% GENERATED FILE. Edit canonical lessons, then run npm run generate:books.
% canonical-source-hash: fnv1a64:73d988572de9bda3
% canonical-lessons: RU-C406-bagazhnik, RU-C406-bamper, RU-C406-fara, RU-C406-sidenye, RU-C406-dvornik

\chapter{Parts of a car}
\label{ch:ru-parts-of-a-car}
\hlchaptermodality{406}

\begin{chapteropening}
\textbf{By the end of this chapter:} \emph{I can name parts of a car: the boot, the bumper, the headlights, the seats and the wipers.}

Complete the last lesson of chapter 406: I can name parts of a car: the boot, the bumper, the headlights, the seats and the wipers.
\end{chapteropening}

\section[\texorpdfstring{\ru{багажник} (bagázhnik)}{багажник (bagázhnik)}]{\ru{багажник} (bagázhnik) — a boot (of a car); a roof rack}

\label{lesson:RU-C406-bagazhnik}



\begin{quote}
Before the new one: say the Russian for to pump up, then the Russian for to spin.
\end{quote}

\subsection*{You'll want to know: \ru{багажник}}
\textbf{\ru{багажник}} (\emph{bagázhnik}) — \textquotedblleft{}a boot (of a car); a roof rack\textquotedblright{}.

\begin{quote}
\ru{Чемоданы} \ru{лежат} \ru{в} \ru{багажнике}.
\end{quote}

\begin{quote}
\emph{Chemodány lezhát v bagázhnike.}
\end{quote}

\begin{quote}
The suitcases are in the boot.
\end{quote}

Built on \emph{\ru{багаж}} — luggage. In American English, the trunk.

\begin{cousinweb}
\emph{\ru{багажник}} is built on \emph{\ru{багаж}} 'luggage' with the ending \emph{-\ru{ник}}, the place for it. \emph{\ru{Багаж}} was borrowed from French \emph{bagage} 'baggage, luggage' (EtymDB); EtymDB does not list \emph{\ru{багажник}} itself.
\end{cousinweb}

\begin{grammarlens}[title={What you've built}]
One more word for telling, explaining and describing at B1.
\end{grammarlens}

\subsection*{Your turn}
\begin{itemize}
\raggedright
  \item \emph{Say it:} \emph{bagázhnik}
  \item \emph{Say it:} \emph{bagázhnik}, once more
  \item \emph{Say it:} a sentence of your own with \emph{bagázhnik}
  \item \emph{Recall:} say the Russian for to pump up, then the Russian for to spin, then say \emph{bagázhnik} again
\end{itemize}

\begin{tcolorbox}[breakable,colback=teal!4,colframe=teal!35!black,title={Before you move on}]
What does \ru{багажник} mean? (\textquotedblleft{}A boot (of a car); a roof rack\textquotedblright{}.) Say it once more.
\end{tcolorbox}
\section[\texorpdfstring{\ru{бампер} (bámper)}{бампер (bámper)}]{\ru{бампер} (bámper) — a bumper}

\label{lesson:RU-C406-bamper}



\begin{quote}
Before the new one: say the Russian for to spin, then the Russian for a boot.
\end{quote}

\subsection*{You'll want to know: \ru{бампер}}
\textbf{\ru{бампер}} (\emph{bámper}) — \textquotedblleft{}a bumper\textquotedblright{}.

\begin{quote}
\ru{После} \ru{аварии} \ru{бампер} \ru{пришлось} \ru{поменять}.
\end{quote}

\begin{quote}
\emph{Pósle avárii bámper prishlós pomenyát.}
\end{quote}

\begin{quote}
After the accident the bumper had to be replaced.
\end{quote}

Masculine; the stress is on the first syllable.

\begin{cousinweb}
\emph{\ru{бампер}} was borrowed from English \emph{bumper} (EtymDB), the bar that takes the bump.
\end{cousinweb}

\begin{grammarlens}[title={What you've built}]
One more word for telling, explaining and describing at B1.
\end{grammarlens}

\subsection*{Your turn}
\begin{itemize}
\raggedright
  \item \emph{Say it:} \emph{bámper}
  \item \emph{Say it:} \emph{bámper}, once more
  \item \emph{Say it:} a sentence of your own with \emph{bámper}
  \item \emph{Recall:} say the Russian for to spin, then the Russian for a boot, then say \emph{bámper} again
\end{itemize}

\begin{tcolorbox}[breakable,colback=teal!4,colframe=teal!35!black,title={Before you move on}]
What does \ru{бампер} mean? (\textquotedblleft{}A bumper\textquotedblright{}.) Say it once more.
\end{tcolorbox}
\section[\texorpdfstring{\ru{фара} (fára)}{фара (fára)}]{\ru{фара} (fára) — a headlight}

\label{lesson:RU-C406-fara}



\begin{quote}
Before the new one: say the Russian for a boot, then the Russian for a bumper.
\end{quote}

\subsection*{You'll want to know: \ru{фара}}
\textbf{\ru{фара}} (\emph{fára}) — \textquotedblleft{}a headlight\textquotedblright{}.

\begin{quote}
\ru{Вечером} \ru{водитель} \ru{включил} \ru{фары}.
\end{quote}

\begin{quote}
\emph{Vécherom vodítel vklyuchíl fáry.}
\end{quote}

\begin{quote}
In the evening the driver switched on the headlights.
\end{quote}

Usually plural: \emph{\ru{фары}}.

\begin{cousinweb}
\emph{\ru{фара}} was borrowed from French \emph{phare} 'a lighthouse; a headlight' (EtymDB). \emph{Phare} goes back to Latin \emph{pharus} 'a lighthouse', from Greek \emph{pháros} (EtymDB).
\end{cousinweb}

\begin{grammarlens}[title={What you've built}]
One more word for telling, explaining and describing at B1.
\end{grammarlens}

\subsection*{Your turn}
\begin{itemize}
\raggedright
  \item \emph{Say it:} \emph{fára}
  \item \emph{Say it:} \emph{fára}, once more
  \item \emph{Say it:} a sentence of your own with \emph{fára}
  \item \emph{Recall:} say the Russian for a boot, then the Russian for a bumper, then say \emph{fára} again
\end{itemize}

\begin{tcolorbox}[breakable,colback=teal!4,colframe=teal!35!black,title={Before you move on}]
What does \ru{фара} mean? (\textquotedblleft{}A headlight\textquotedblright{}.) Say it once more.
\end{tcolorbox}
\section[\texorpdfstring{\ru{сиденье} (sidénye)}{сиденье (sidénye)}]{\ru{сиденье} (sidénye) — a seat (in a car or a bus)}

\label{lesson:RU-C406-sidenye}



\begin{quote}
Before the new one: say the Russian for a bumper, then the Russian for a headlight.
\end{quote}

\subsection*{You'll want to know: \ru{сиденье}}
\textbf{\ru{сиденье}} (\emph{sidénye}) — \textquotedblleft{}a seat (in a car or a bus)\textquotedblright{}.

\begin{quote}
\ru{Дети} \ru{сидят} \ru{на} \ru{заднем} \ru{сиденье}.
\end{quote}

\begin{quote}
\emph{Déti sidyát na zádnem sidénye.}
\end{quote}

\begin{quote}
The children are sitting on the back seat.
\end{quote}

Neuter. Also spelled \emph{\ru{сидение}}.

\begin{cousinweb}
\emph{\ru{сиденье}} is built on \emph{\ru{сидеть}} 'to sit', which is inherited from Proto-Slavic \emph{sěděti} 'to sit' (EtymDB).
\end{cousinweb}

\begin{grammarlens}[title={What you've built}]
One more word for telling, explaining and describing at B1.
\end{grammarlens}

\subsection*{Your turn}
\begin{itemize}
\raggedright
  \item \emph{Say it:} \emph{sidénye}
  \item \emph{Say it:} \emph{sidénye}, once more
  \item \emph{Say it:} a sentence of your own with \emph{sidénye}
  \item \emph{Recall:} say the Russian for a bumper, then the Russian for a headlight, then say \emph{sidénye} again
\end{itemize}

\begin{tcolorbox}[breakable,colback=teal!4,colframe=teal!35!black,title={Before you move on}]
What does \ru{сиденье} mean? (\textquotedblleft{}A seat (in a car or a bus)\textquotedblright{}.) Say it once more.
\end{tcolorbox}
\section[\texorpdfstring{\ru{дворник} (dvórnik)}{дворник (dvórnik)}]{\ru{дворник} (dvórnik) — a windscreen wiper; a caretaker, a street sweeper}

\label{lesson:RU-C406-dvornik}



\begin{quote}
Before the new one: say the Russian for a headlight, then the Russian for a seat.
\end{quote}

\subsection*{You'll want to know: \ru{дворник}}
\textbf{\ru{дворник}} (\emph{dvórnik}) — \textquotedblleft{}a windscreen wiper; a caretaker, a street sweeper\textquotedblright{}.

\begin{quote}
\ru{Шёл} \ru{дождь}, \ru{и} \ru{дворники} \ru{работали} \ru{всю} \ru{дорогу}.
\end{quote}

\begin{quote}
\emph{Shyol dozhd, i dvórniki rabótali vsyu dorógu.}
\end{quote}

\begin{quote}
It was raining, and the wipers worked all the way.
\end{quote}

The first meaning is a person who sweeps a yard; the wipers sweep the glass.

\begin{cousinweb}
\emph{\ru{дворник}} is built on \emph{\ru{двор}} 'a yard' (EtymDB), which is inherited from Proto-Slavic \emph{dvorŭ} 'a court'. EtymDB glosses \emph{\ru{дворник}} as the one who keeps the yard and the pavement; the wiper that sweeps the glass took the name later.
\end{cousinweb}

\begin{grammarlens}[title={What you've built}]
One more word for telling, explaining and describing at B1.
\end{grammarlens}

\subsection*{Your turn}
\begin{itemize}
\raggedright
  \item \emph{Say it:} \emph{dvórnik}
  \item \emph{Say it:} \emph{dvórnik}, once more
  \item \emph{Say it:} a sentence of your own with \emph{dvórnik}
  \item \emph{Recall:} say the Russian for a headlight, then the Russian for a seat, then say \emph{dvórnik} again
\end{itemize}

\begin{tcolorbox}[breakable,colback=teal!4,colframe=teal!35!black,title={Before you move on}]
What does \ru{дворник} mean? (\textquotedblleft{}A windscreen wiper; a caretaker, a street sweeper\textquotedblright{}.) Say it once more.
\end{tcolorbox}
